\chapter{系统分析与总体设计}

\section{需求分析}

系统用户分为学生用户、教师/评委和管理员三类。学生用户关注上传、分类、预览、转换和 PDF 编辑；教师/评委关注系统是否具备真实功能和可演示证据；管理员关注任务队列、失败任务、系统健康和数据驾驶舱。

\section{功能模块设计}

\begin{table}[htbp]
  \centering
  \caption{系统功能模块划分}
  \label{tab:function-modules}
  \begin{tabular}{lll}
    \toprule
    模块 & 主要功能 & 数据来源 \\
    \midrule
    个人首页 & 文件数、空间、转换成功、失败任务 & 文件 API、任务 API \\
    我的云盘 & 上传、分类、预览、下载、重命名、删除 & MinIO/HDFS \\
    转换中心 & Office/PDF/图片/Markdown 转换 & 转换 API、任务队列 \\
    PDF 编辑室 & 预览、批注、页面整理、打印导出 & PDF 页面提取 API \\
    管理驾驶舱 & 遥测分析、质量指标、拓扑与状态 & Spark 聚合数据 \\
    系统状态 & 服务健康、队列长度、近期失败 & health/task/log API \\
    \bottomrule
  \end{tabular}
\end{table}

\section{数据流设计}

用户上传文件后，前端调用文件上传 API，将文件写入对象存储。预览模块根据文件类型生成 PDF、图片、文本或 HTML 预览内容。转换中心提交异步任务，任务完成后生成转换产物并写回云盘。PDF 编辑室加载 PDF 页面图，前端记录批注对象；页面重排导出由后端根据页面配置生成新 PDF，批注打印导出由浏览器打印渲染完成。

\section{合规与安全设计}

系统采用以下安全设计：上传和预览接口分离，避免预览缓存覆盖原文件；文件删除和重命名由用户主动触发；Dashboard 在接口失败时展示不可用状态而不是假数据；通知面板只展示真实失败任务；日志仅用于实验追踪和异常定位。

\section{数据库与对象结构}

系统文件实体包含 \texttt{object\_name}、\texttt{filename}、\texttt{size}、\texttt{last\_modified} 和 \texttt{content\_type}。任务实体包含 \texttt{task\_id}、\texttt{kind}、\texttt{status}、\texttt{created\_at}、\texttt{completed\_at}、\texttt{error} 和 \texttt{result\_url}。这些字段支撑首页真实指标、任务列表、失败通知和文档测试表。

